\part{Quantum search}\label{part:search}

\begin{frame}
  \partpage
\end{frame}

%%%%%%%%%%%%%%%%%%%%%%%%%%%%%%%%%%%%%%%%%%%%%%%%%%%%%%%%%%%%%%%%%%%%%%%%%%%%%%%%

%% Quantum search

%%%%%%%%%%%%%%%%%%%%%%%%%%%%%%%%%%%%%%%%%%%%%%%%%%%%%%%%%%%%%%%%%%%%%%%%%%%%%%%%

\begin{frame}{Unstructured search}

Quantum computers can quadratically outperform classical computers at a very basic computational 
task, called unstructured search.

\bigskip
There is a set $X$ containing $N$ items, some of which are marked

\bigskip
We are given a Boolean black box $f : X \rightarrow \{0,1\}$ that indicates whether a given item is marked

\bigskip
The problem is to decide if any item is marked, or alternatively, to find a marked item given that one exists

\end{frame}

%%%%%%%%%%%%%%%%%%%%%%%%%%%%%%%%%%%%%%%%%%%%%%%%%%%%%%%%%%%%%%%%%%%%%%%%%%%%%%%%%

%% Applications of unstructured search

%%%%%%%%%%%%%%%%%%%%%%%%%%%%%%%%%%%%%%%%%%%%%%%%%%%%%%%%%%%%%%%%%%%%%%%%%%%%%%%%%

\begin{frame}{Applications of unstructured search}

Unstructured search can be thought of as a model for solving problems in NP by brute force search

\bigskip
If a problem is in NP, then we can efficiently recognize a solution, so one way to find a solution is to solve unstructured search

\bigskip
Of course, this may not be the best way to find a solution in general, even if the problem is NP-hard

\bigskip
We don't know if NP-hard problems are really ``unstructured''
\end{frame}

%%%%%%%%%%%%%%%%%%%%%%%%%%%%%%%%%%%%%%%%%%%%%%%%%%%%%%%%%%%%%%%%%%%%%%%%%%%%%%%%%

%% Complexity of unstructured search

%%%%%%%%%%%%%%%%%%%%%%%%%%%%%%%%%%%%%%%%%%%%%%%%%%%%%%%%%%%%%%%%%%%%%%%%%%%%%%%%%

\begin{frame}{Unstructured search}

It is obvious that even a randomized classical algorithm needs $\Omega(N)$ queries to decide if any item is marked  

\bigskip
On the other hand, a quantum algorithm can do much better!
\end{frame}

%%%%%%%%%%%%%%%%%%%%%%%%%%%%%%%%%%%%%%%%%%%%%%%%%%%%%%%%%%%%%%%%%%%%%%%%%%%%%%%%%

%% Phase oracle

%%%%%%%%%%%%%%%%%%%%%%%%%%%%%%%%%%%%%%%%%%%%%%%%%%%%%%%%%%%%%%%%%%%%%%%%%%%%%%%%%

\begin{frame}{Phase oracle}

We assume that we a unitary operator $U$ satisfying

\[
  U|x\> = (-1)^{f(x)} |x\> = 
  \begin{cases}
  	\phantom{-} |x\> & x \text{~is not marked} \\
             -  |x\> & x \text{~is marked}
  \end{cases}
\]

\end{frame}

%%%%%%%%%%%%%%%%%%%%%%%%%%%%%%%%%%%%%%%%%%%%%%%%%%%%%%%%%%%%%%%%%%%%%%%%%%%%%%%%%

%% Target state

%%%%%%%%%%%%%%%%%%%%%%%%%%%%%%%%%%%%%%%%%%%%%%%%%%%%%%%%%%%%%%%%%%%%%%%%%%%%%%%%%

\begin{frame}{Target state}

We consider the case where there is exactly one $x\in X$ element that is marked; call this element $m$

\bigskip
Our goal is to prepare the state $|m\>$

\end{frame}

%%%%%%%%%%%%%%%%%%%%%%%%%%%%%%%%%%%%%%%%%%%%%%%%%%%%%%%%%%%%%%%%%%%%%%%%%%%%%%%%%

%% Initial state

%%%%%%%%%%%%%%%%%%%%%%%%%%%%%%%%%%%%%%%%%%%%%%%%%%%%%%%%%%%%%%%%%%%%%%%%%%%%%%%%%

\begin{frame}{Initial state}

We have no information about which item might be marked

\bigskip
$\Rightarrow$ We take

\[
|\psi\> := \frac{1}{\sqrt N} \sum_{x=1}^N |x\>
\]

as the initial state

\end{frame}

%%%%%%%%%%%%%%%%%%%%%%%%%%%%%%%%%%%%%%%%%%%%%%%%%%%%%%%%%%%%%%%%%%%%%%%%%%%%%%%%%

%% Idea behind Grover search

%%%%%%%%%%%%%%%%%%%%%%%%%%%%%%%%%%%%%%%%%%%%%%%%%%%%%%%%%%%%%%%%%%%%%%%%%%%%%%%%%


\begin{frame}{Rough idea behind Grover search}

We start with the initial state $|\psi\>$

\bigskip
We prepare the target state $|m\>$ by implementing a rotation that moves $|\psi\>$ toward $|m\>$

\bigskip
We realize the rotation with the help of two reflections

\end{frame}

%%%%%%%%%%%%%%%%%%%%%%%%%%%%%%%%%%%%%%%%%%%%%%%%%%%%%%%%%%%%%%%%%%%%%%%%%%%%%%%%%

%% Reflection -- visualization

%%%%%%%%%%%%%%%%%%%%%%%%%%%%%%%%%%%%%%%%%%%%%%%%%%%%%%%%%%%%%%%%%%%%%%%%%%%%%%%%%

\begin{frame}{Visualization of a reflection in $\R^2$}

\begin{center}
\scalebox{.65}{
\multiinclude[<+-| handout:+->][format=pdf]{G1}
}
\end{center}

\end{frame}

%%%%%%%%%%%%%%%%%%%%%%%%%%%%%%%%%%%%%%%%%%%%%%%%%%%%%%%%%%%%%%%%%%%%%%%%%%%%%%%%%

%% Reflections

%%%%%%%%%%%%%%%%%%%%%%%%%%%%%%%%%%%%%%%%%%%%%%%%%%%%%%%%%%%%%%%%%%%%%%%%%%%%%%%%%

\begin{frame}{Reflections}

$U=I-2|m\>\<m|$ is a reflection about the target state $|m\>$

\bigskip
$V=I-2|\psi\>\<\psi|$ is the reflection around about the initial state $|\psi\>$:

\begin{eqnarray*}
V |\psi\>       & = & -   \, |\psi\> \\
V |\psi^\perp\> & = &  \quad |\psi^\perp\>
\end{eqnarray*}

\bigskip
for any state $|\psi^\perp\>$ orthogonal to $|\psi\>$

\end{frame}

%%%%%%%%%%%%%%%%%%%%%%%%%%%%%%%%%%%%%%%%%%%%%%%%%%%%%%%%%%%%%%%%%%%%%%%%%%%%%%%%%

%% Structure of Grover search

%%%%%%%%%%%%%%%%%%%%%%%%%%%%%%%%%%%%%%%%%%%%%%%%%%%%%%%%%%%%%%%%%%%%%%%%%%%%%%%%%

\begin{frame}{Structure of Grover}

\bigskip
The algorithm is as follows: 

\bigskip
\begin{itemize}
\item start in $|\psi\>$, 

\bigskip
\item apply the Grover iteration $G := V\, U$ some number of times, 

\bigskip
\item make a measurement, and hope that the outcome is $m$
\end{itemize}
\end{frame}

%%%%%%%%%%%%%%%%%%%%%%%%%%%%%%%%%%%%%%%%%%%%%%%%%%%%%%%%%%%%%%%%%%%%%%%%%%%%%%%%%

%% Invariant subspace

%%%%%%%%%%%%%%%%%%%%%%%%%%%%%%%%%%%%%%%%%%%%%%%%%%%%%%%%%%%%%%%%%%%%%%%%%%%%%%%%%

\begin{frame}{Invariant subspace}

Observe that $\spn\{|m\>,|\psi\>\}$ is a $U$- and $V$-invariant subspace, and both the inital and target states belong to 
this subspace

\bigskip
$\Rightarrow$ It suffices to understand the restriction of $VU$ to this subspace

\bigskip
Consider an orthonormal basis $\{|m\>,|\phi\>\}$ for $\spn\{|m\>,|\psi\>\}$

\bigskip
The Gram-Schmidt process yields

\bigskip
\[
  |\phi\> = \frac{|\psi\> - \alpha |m\>}{\sqrt{1-\alpha^2}}
\]

\bigskip
where $\alpha := \<m|\psi\> = 1/\sqrt{N}$
\end{frame}

%%%%%%%%%%%%%%%%%%%%%%%%%%%%%%%%%%%%%%%%%%%%%%%%%%%%%%%%%%%%%%%%%%%%%%%%%%%%%%%%%

%% Invariant subspace

%%%%%%%%%%%%%%%%%%%%%%%%%%%%%%%%%%%%%%%%%%%%%%%%%%%%%%%%%%%%%%%%%%%%%%%%%%%%%%%%%

\begin{frame}{Invariant subspace}

Now in the basis $\{|m\>,|\phi\>\}$, we have

\begin{eqnarray*}  
|\psi\> 
& = & 
\sin\theta |m\> + \cos\theta |\phi\> \mbox{ where } \sin\theta = \<m|\psi\> = 1/\sqrt{N} \\
&   & \\
U & = & 
\begin{pmatrix}-1 & 0 \\ 0 & 1\end{pmatrix} \\
&   & \\ \pause
V 
& = & 
I-2|\psi\>\<\psi| \\ \pause
& = &   
\begin{pmatrix} 1 & 0 \\ 0 & 1 \end{pmatrix} 
-2 
\begin{pmatrix} \sin\theta \\ \cos\theta \end{pmatrix} 
\begin{pmatrix} \sin\theta & \cos\theta\end{pmatrix} \\ \pause
& = & 
\begin{pmatrix} 
  1-2\sin^2\theta         & -2\sin\theta\cos\theta \\ 
   -2\sin\theta\cos\theta & 1-2\cos^2\theta
\end{pmatrix} \\ \pause
& = &
-
\begin{pmatrix} 
  -\cos 2\theta & \sin 2\theta \\
   \sin 2\theta & \cos 2 \theta 
\end{pmatrix} 
\end{eqnarray*}
\end{frame}

%%%%%%%%%%%%%%%%%%%%%%%%%%%%%%%%%%%%%%%%%%%%%%%%%%%%%%%%%%%%%%%%%%%%%%%%%%%%%%%%%

%% Grover iteration 

%%%%%%%%%%%%%%%%%%%%%%%%%%%%%%%%%%%%%%%%%%%%%%%%%%%%%%%%%%%%%%%%%%%%%%%%%%%%%%%%%

\begin{frame}{Grover iteration within the invariant subspace}

$\Rightarrow$ We find

\begin{eqnarray*}
V \, U 
& = &
-
\begin{pmatrix} 
  -\cos 2\theta & \sin 2\theta \\
   \sin 2\theta & \cos 2 \theta 
\end{pmatrix} 
\,
\begin{pmatrix} -1 & 0 \\ 0 & 1 \end{pmatrix} \\ \pause
& = &
- 
\begin{pmatrix}  
	  \cos 2\theta & \sin 2\theta \\
  - \sin 2\theta & \cos 2\theta
\end{pmatrix}
\end{eqnarray*}

\bigskip
This is a rotation up to a minus sign
\end{frame}

%%%%%%%%%%%%%%%%%%%%%%%%%%%%%%%%%%%%%%%%%%%%%%%%%%%%%%%%%%%%%%%%%%%%%%%%%%%%%%%%%

%% first Grover iteration

%%%%%%%%%%%%%%%%%%%%%%%%%%%%%%%%%%%%%%%%%%%%%%%%%%%%%%%%%%%%%%%%%%%%%%%%%%%%%%%%%

\begin{frame}{Visualization of first Grover iteration}

\begin{center}
\scalebox{.48}{
\multiinclude[<+-| handout:+->][format=pdf]{G2}
}
\end{center}
\end{frame}

%%%%%%%%%%%%%%%%%%%%%%%%%%%%%%%%%%%%%%%%%%%%%%%%%%%%%%%%%%%%%%%%%%%%%%%%%%%%%%%%%

%% Grover search

%%%%%%%%%%%%%%%%%%%%%%%%%%%%%%%%%%%%%%%%%%%%%%%%%%%%%%%%%%%%%%%%%%%%%%%%%%%%%%%%%

\begin{frame}{Grover search}

Geometrically, $U$ is a reflection around the $|m\>$ axis and $V$ is a reflection around the $|\psi\>$ axis, 
which is almost but not quite orthogonal to the $|m\>$ axis

\bigskip
The product of these two reflections is a clockwise rotation by an angle $2\theta$, up to an overall minus sign

\bigskip
From this geometric picture, or by explicit calculation using trig identities, it is easy to verify that

\bigskip
\[
  (VU)^k = (-1)^k \begin{pmatrix} \cos2k\theta & \sin2k\theta \\
                       - \sin2k\theta & \cos2k\theta\end{pmatrix}
\]

\end{frame}

%%%%%%%%%%%%%%%%%%%%%%%%%%%%%%%%%%%%%%%%%%%%%%%%%%%%%%%%%%%%%%%%%%%%%%%%%%%%%%%%%

%% Grover search

%%%%%%%%%%%%%%%%%%%%%%%%%%%%%%%%%%%%%%%%%%%%%%%%%%%%%%%%%%%%%%%%%%%%%%%%%%%%%%%%%

\begin{frame}{Grover search}

Recall that our initial state is $|\psi\> = \sin\theta|m\> + \cos\theta|\phi\>$

\bigskip
How large should $k$ be before $(VU)^k|\psi\>$ is close to $|m\>$?  

\bigskip
We start an angle $\theta$ from the $|\phi\>$ axis and rotate 
toward $|m\>$ by an angle $2\theta$ per iteration

\bigskip
$\Rightarrow$ To rotate by $\pi/2$, we 
need 

\begin{eqnarray*}
\theta + 2k\theta & =       & \pi/2 \quad \\ & & \\
k                 & \approx & \frac{\pi}{4}\theta^{-1} \approx \frac{\pi}{4} \sqrt{N}
\end{eqnarray*}

\end{frame}

%%%%%%%%%%%%%%%%%%%%%%%%%%%%%%%%%%%%%%%%%%%%%%%%%%%%%%%%%%%%%%%%%%%%%%%%%%%%%%%%%

%% Grover search

%%%%%%%%%%%%%%%%%%%%%%%%%%%%%%%%%%%%%%%%%%%%%%%%%%%%%%%%%%%%%%%%%%%%%%%%%%%%%%%%%

\begin{frame}{Grover search}

It is easy to calculate that 

\medskip
\[
|\<m|(VU)^k|\psi\>|^2 = \sin^2((2k+1)\theta)
\]

\bigskip
This is the probability that, after $k$ steps of the algorithm, a measurement reveals the marked state

\bigskip
We are solving a completely unstructured search problem with $N$ possible solutions, yet we 
can find a unique solution in only $O(\sqrt N)$ queries!  

\bigskip
While this is only a polynomial separation, it is very generic, and it is surprising that we can obtain a speedup for a search in which we have so little 
information to go on

\end{frame}

%%%%%%%%%%%%%%%%%%%%%%%%%%%%%%%%%%%%%%%%%%%%%%%%%%%%%%%%%%%%%%%%%%%%%%%%%%%%%%%%%

\begin{frame}{Grover search}

It can also be shown that this quantum algorithm is optimal

\bigskip
Any quantum algorithm needs at least $\Omega(\sqrt N)$ queries to find a marked item (or even to decide if some item is marked)

\end{frame}

%%%%%%%%%%%%%%%%%%%%%%%%%%%%%%%%%%%%%%%%%%%%%%%%%%%%%%%%%%%%%%%%%%%%%%%%%%%%%%%%%%

\begin{frame}{Multiple solutions}

Assume that there are $t$ marked items

\bigskip
$\Rightarrow$ There is a two-dimensional invariant subspace spanned by $\spn\{|\mu\>,|\psi\>\}$ where

\[
|\mu\> = \frac{1}{\sqrt{t}} \sum_{x \, \mathsf{marked}} |x\>
\]

is the uniform superposition of all solutions

\bigskip
The Gram-Schmidt process yields the ONB $\{|\mu\>,|\phi\>\}$ where

\[
|\phi\> = \frac{1}{\sqrt{N-t}} \sum_{x \, \mathsf{unmarked}} |x\>
\]

is the uniform superposition of all non-solutions

\end{frame}

%%%%%%%%%%%%%%%%%%%%%%%%%%%%%%%%%%%%%%%%%%%%%%%%%%%%%%%%%%%%%%%%%%%%%%%%%%%%%%%%%

%% Invariant subspace

%%%%%%%%%%%%%%%%%%%%%%%%%%%%%%%%%%%%%%%%%%%%%%%%%%%%%%%%%%%%%%%%%%%%%%%%%%%%%%%%%

\begin{frame}{Invariant subspace}

Now in the basis $\{|\mu\>,|\phi\>\}$, we have

\begin{eqnarray*}  
	|\psi\> & = & \sin\theta |\mu\> + \cos\theta |\phi\> \mbox{ where } \sin\theta = \<\mu|\psi\> = \sqrt{\frac{t}{N}} \\
					&   & \\
  V U & = & - \begin{pmatrix} \cos2\theta   & \sin2\theta \\
                              - \sin2\theta & \cos2\theta
              \end{pmatrix}
\end{eqnarray*}
\end{frame}

%%%%%%%%%%%%%%%%%%%%%%%%%%%%%%%%%%%%%%%%%%%%%%%%%%%%%%%%%%%%%%%%%%%%%%%%%%%%%%%%%%

\begin{frame}{Overshooting}

The success probability is given by

\[
\sin((2k + 1)\theta) \mbox{ where } \sin\theta = \sqrt{\frac{t}{N}}
\]

\bigskip
$\Rightarrow$ We need to apply $VU$

\[
k \approx \frac{\pi}{4} \sqrt{\frac{N}{t}}
\]

times

\bigskip
Due to the oscillatory behaviour of the success probability it is important not to overshoot, i.e., to choose a number of iterations that it too large, so that the probability starts decreasing
\end{frame}

%%%%%%%%%%%%%%%%%%%%%%%%%%%%%%%%%%%%%%%%%%%%%%%%%%%%%%%%%%%%%%%%%%%%%%%%%%%%%%%%%%

\begin{frame}{Quantum counting}

The eigenvalues of 

\[
- V U = \begin{pmatrix} \cos2\theta   & \sin2\theta \\
                              - \sin2\theta & \cos2\theta
              \end{pmatrix}
\]


are $e^{i 2\theta}$ and $e^{-i 2\theta}$

\bigskip
The initial state $|\psi\>$ is a superposition of the two eigenvectors corresponding to the above two eigenvalues

\bigskip
$\Rightarrow$ Using phase estimation, we can obtain an estimate $\tilde{\theta}$ such that

\[
|\theta - \tilde{\theta} | \le \epsilon
\]

\bigskip
by invoking the controlled version of $-VU$ 

\[
O(1/\epsilon) \mbox{ times}
\]

\end{frame}

%%%%%%%%%%%%%%%%%%%%%%%%%%%%%%%%%%%%%%%%%%%%%%%%%%%%%%%%%%%%%%%%%%%%%%%%%%%%%%%%%%

\begin{frame}{Quantum counting}

Using the estimate $\tilde{\theta}$, we obtain an estimate $\tilde{t}$ satifying

\[
|t - \tilde{t} | \le (2\sqrt{t N} + \epsilon) \, \epsilon
\]

\end{frame}

%%%%%%%%%%%%%%%%%%%%%%%%%%%%%%%%%%%%%%%%%%%%%%%%%%%%%%%%%%%%%%%%%%%%%%%%%%%%%%%%%%

\begin{frame}{Quantum counting}

We use the following two inequalities:
\begin{eqnarray*}
|\sin\theta + \sin\tilde{\theta}| & \le & 2 \sin\theta + |\theta - \tilde{\theta} |  \, \le \, 2\sqrt{\frac{t}{N}} + \epsilon\\
|\sin\theta - \sin\tilde{\theta}| & \le & |\theta - \tilde{\theta}|  \, \le \, \epsilon
\end{eqnarray*}

\bigskip
We have
\begin{eqnarray*}
\left| \frac{t}{N} - \frac{\tilde{t}}{N} \right| 
& = &
|\sin^2 \theta - \sin^2\tilde{\theta}| \\
& = &
|\sin\theta + \sin\tilde{\theta}| \, |\sin\theta - \sin\tilde{\theta}| \\
& \le &
\left(
2 \sqrt{\frac{t}{N}} + \epsilon
\right) \, \epsilon
\end{eqnarray*}

\end{frame}

%%%%%%%%%%%%%%%%%%%%%%%%%%%%%%%%%%%%%%%%%%%%%%%%%%%%%%%%%%%%%%%%%%%%%%%%%%%%%%%%%%

\begin{frame}{Amplitude amplification}

Assume that there is a classical (randomized) algorithm that produces a solution to some problem with probability $p$

\bigskip
Assume that we can recognized if the output produced by the algorithm is a valid solution or not

\bigskip
$\Rightarrow$ We repeat the algorithm until we obtain a solution

\bigskip
The expected number of times we have to repeat is $O(1/p)$ (geometric random variable)

\bigskip
Quantum amplitude amplification makes it possible to reduce the complexity to $O(1/\sqrt{p})$

\end{frame}

%%%%%%%%%%%%%%%%%%%%%%%%%%%%%%%%%%%%%%%%%%%%%%%%%%%%%%%%%%%%%%%%%%%%%%%%%%%%%%%%%%